\section{Fazit und Ausblick}\label{sec:fazit}
Die Forschungsfrage lässt sich verdichtet beantworten. Die Potenziale liegen in beschleunigter Erstellung, besserer Dokumentation und niedrigeren Einstiegshürden, die Risiken in unsicherem Code, Overreliance, nachlassender Kompetenz und verzerrter Selbsteinschätzung. Beide Seiten treffen dieselben Praktiken. Ob Nutzen oder Schaden überwiegt, entscheidet sich daher nicht an der Werkzeugwahl, sondern daran, ob Review, Definition of Done und Retrospektive an generierten Code angepasst werden; die Empfehlungen aus Kapitel~\ref{sec:wuerdigung} setzen dort an.

Offen bleiben drei Fragen. Erstens die Langzeitwirkung auf die Kompetenzentwicklung, da die vorliegenden Untersuchungen kurze Zeiträume abbilden. Zweitens die Wirkung \glslink{agentic}{agentischer Werkzeuge}: Sie ändern nach Kapitel~\ref{subsec:generatoren} selbstständig Dateien, sodass die menschliche Prüfung erst nach mehreren Arbeitsschritten ansetzt. Drittens, ob sich die Rolle der Entwicklung dauerhaft vom Ausführen zum Überwachen verschiebt -- eine Verschiebung, die für frühere Automatisierungswellen bereits beschrieben ist.\footcite[\vglf][\pagef 1]{Shen.2026}
